\section{Solving the status constraints}
\label{sec:status-span}

We now construct the options forbidden by \Cref{prop:limit-overlap}.
The argument takes place in the common signature model, but the
no-prediction alternative will require a transfer back to global
functions on the finite key spaces.
We show that a positive mass of permitted pairs with soluble scalar
constraints forces positive key overlap. Applied to the prepared
coloured law, the collision criterion then supplies a four-hole event
between different colours.

Write the endpoints of the two units as $(A,1),(A,2)$ and
$(B,1),(B,2)$. In this section $h_{A,i}$ or $h_{B,j}$ denotes an atom's
role bit $a(\mathrm{atom})$, not the channel dimension. The basis status
vector at an endpoint is denoted by $\mathbf x$, and its dimension is
the recorded dimension of that endpoint's effective space $C_i$.

\subsection{Generators and the coupled linear system}

Fix two orientation arrays $\theta_A,\theta_B$ satisfying the
almost-sure conclusions of \Cref{lem:limit-model}. Some labels at each
endpoint may be excluded. In particular, the mandatory support
exclusions from \Cref{eq:label-budget} are always excluded.

\begin{definition}[Available generator]
\label{def:status-generator}
On the cross edge $(i,j)$, a vector
\begin{equation}
\label{eq:status-generator-vector}
g=(q,h_{A,i},h_{B,j},\mathbf x_{A,i},\mathbf x_{B,j},p_{A,i},p_{B,j})
\end{equation}
is an available generator if some common tag, surviving label at each
end, and allowed flavors give positive overlap as in
\eqref{eq:limit-positive-generator} for these entries. Entries of the
tester summary not displayed in \eqref{eq:status-generator-vector}
are summed over. A generator is retained with one such certificate of
its tags, labels, and flavors.

Let $V_{ij}$ be the binary span of these generators. Its \emph{base
projection} forgets the final two $p$ entries and has target
\[
B_{ij}=\F\oplus\F\oplus\F\oplus
             \F^{\dim C_{A,i}}\oplus\F^{\dim C_{B,j}}.
\]
\end{definition}

For any surviving pair of labels and any tag, the full generic flavor
gives at least one available generator above every base vector. Indeed, at fixed $q$
the sum of the status densities with any requested role/basis assignment
is bounded below almost everywhere by
\Cref{lem:limit-model}. The product of the two such sums has positive
integral on the common signature space. Splitting its finite sum over
the two $p$ bits yields an available generator. Both diagonal bits are
feasible in the full generic flavor. Thus
\begin{equation}
\label{eq:status-full-projection}
V_{ij}\longrightarrow B_{ij}\text{ is surjective},\qquad
\dim\ker(V_{ij}\longrightarrow B_{ij})\le2.
\end{equation}
This remains true after any exclusions leaving surviving generic labels.

Fix a small table. Its local target basis vectors are the evaluations of
\[
(a+(u_{B,j}U_{A,i})^*)|_{C_{A,i}},\qquad
(a+(u_{A,i}U_{B,j})^*)|_{C_{B,j}},
\]
in the recorded bases. Denote these vectors by
$\mathbf t_{A,i;j}$ and $\mathbf t_{B,j;i}$, respectively. We seek one
sum vector $v_{ij}\in V_{ij}$ for each cross edge, subject to
\begin{equation}
\label{eq:status-local-targets}
\begin{aligned}
q_{ij}&=1,& h_{A,i;ij}+h_{B,j;ij}&=1,\\
\mathbf x_{A,i;ij}&=\mathbf t_{A,i;j},&
\mathbf x_{B,j;ij}&=\mathbf t_{B,j;i},
\end{aligned}
\end{equation}
\begin{equation}
\label{eq:status-coupled-targets}
\begin{aligned}
\sum_{i=1}^2(p_{A,i;ij}+h_{A,i;ij})&=0 &&(j=1,2),\\
\sum_{j=1}^2(p_{B,j;ij}+h_{B,j;ij})&=0 &&(i=1,2).
\end{aligned}
\end{equation}
These are precisely the scalar recipe equations
\eqref{eq:recipe-scalars}. The occurrence of a given endpoint in two
different sums is distinguished by its edge subscript.
By \Cref{thm:gradient-realization}, these are the effective-space
conditions needed for the full gradients; the attack-atom conditions
are supplied by \eqref{eq:recipe-gradients}. This use does not require
pins to be supported at individual endpoints. The realization theorem
already treats the mixed pin spaces through their projections and the
injection hypothesis.

\begin{lemma}[Dual relations]
\label{lem:status-duality}
Any obstruction to the scalar system has a nonzero vector of weights
$(s_1,s_2,r_1,r_2)\in\F^4$. On every edge span $V_{ij}$ these give
a relation
\begin{equation}
\label{eq:status-dual-relation}
s_j(p_{A,i}+h_{A,i})+r_i(p_{B,j}+h_{B,j})
=\mathbf c_{ij}\cdot\mathbf x_{A,i}
 +\mathbf c'_{ij}\cdot\mathbf x_{B,j}
 +\beta_{ij}(h_{A,i}+h_{B,j})+\gamma_{ij}q
\end{equation}
for suitable coefficients. The scalar system is soluble if, for every
such relation, the sum of the right sides evaluated at the local targets
is zero.
\end{lemma}

\begin{proof}
Apply finite-dimensional duality to the linear constraint map on
$\bigoplus_{i,j}V_{ij}$. The coefficients of the four coupled equations
are $s_j,r_i$, giving \eqref{eq:status-dual-relation} separately on each
edge. If they were all zero, the local projection onto
$(q,h_{A,i}+h_{B,j},\mathbf x_{A,i},\mathbf x_{B,j})$ would be
surjective by \eqref{eq:status-full-projection}, forcing all other dual
coefficients to vanish. Hence a nonzero obstruction uses some coupled
equation. Conversely the constraint target lies in the image exactly
when it annihilates every dual relation. The coupled targets are zero,
so their evaluated contribution is the sum stated in the lemma.
\end{proof}

\subsection{A dual pattern on positive mass gives a global prediction}

Call a pair of arrays \emph{obstructable} if there are total exclusion
sets of size at most $B_*$ at each endpoint, containing the mandatory
exclusions, and some nonzero weights for which
\eqref{eq:status-dual-relation} holds throughout. We do not require the
relation to take a nonzero value on a particular table target. This
stronger event is useful because its complement works for every table.

\begin{lemma}[Extraction of a global prediction]
\label{lem:status-global-prediction}
On the no-prediction alternative of \Cref{prop:status-preparation}, the
probability that the iid base arrays are obstructable is less than $.1$.
\end{lemma}

\begin{proof}
Suppose otherwise. Let $E_s$ be the event that a witnessing relation
has some $s_j=1$, and define $E_r$ similarly. Interchanging the two iid
units interchanges these events, so they have the same probability.
Their union is the obstructable event. Therefore
$\PP(E_s)\ge .05$.

Fix a witness with $s_j=1$. For each first-side endpoint $i$, rewrite
its edge relation as
\begin{align}
p_{A,i}+\mathbf c_{ij}\cdot\mathbf x_{A,i}
                 +(1+\beta_{ij})h_{A,i}
={}&r_i p_{B,j}+\mathbf c'_{ij}\cdot\mathbf x_{B,j}
                 +(r_i+\beta_{ij})h_{B,j}+\gamma_{ij}q.
\label{eq:status-deterministic-comparison}
\end{align}
For any surviving label/flavor choices and any common feasible $(l,q)$,
this equality holds almost surely under the product of the two status
distributions conditional on the same signature. Otherwise a violating
pair of numerical statuses would have positive overlap and would be a
generator violating the relation. The two bits in
\eqref{eq:status-deterministic-comparison} are conditionally independent
at a fixed signature. If two independent bits agree almost surely, both
are constant almost surely with their common value: writing their
success probabilities as $a,b$, the mismatch probability
$a(1-b)+(1-a)b$ is zero only for $(a,b)=(0,0)$ or $(1,1)$.

Compare all surviving first-side choices to one fixed generic surviving
choice at $(B,j)$. Its total conditional status mass is one on every
signature space, and its generic flavor admits either $q$. It follows
that, simultaneously over all tags, feasible bits, labels and flavors
remaining at $(A,i)$,
\begin{equation}
\label{eq:status-extracted-prediction}
p_{A,i}=T(\lambda_{A,i},\mathrm{atom})+\alpha_{A,i}h_{A,i}
                    +f_i(l,q,\xi),
\end{equation}
where $\lambda_{A,i}$ is the basis combination specified by
$\mathbf c_{ij}$ and $\alpha_{A,i}=1+\beta_{ij}$. The coefficients
are common to these choices. The offset is also common across them;
one does not choose a new offset for every label or flavor.

There are very few possible offsets after the second-side array is
fixed. If $r_i=0$, generic positivity of the role and basis at $(B,j)$
forces $\mathbf c'_{ij}=0$ and $\beta_{ij}=0$ in
\eqref{eq:status-deterministic-comparison}. The offset is then either
zero or $q$. If $r_i=1$, it is, up to a common multiple of $q$, the
deterministic value of
\begin{equation}
\label{eq:status-opposite-offset}
p_{B,j}+\mathbf c'\cdot\mathbf x_{B,j}+\alpha' h_{B,j}.
\end{equation}
For this fixed array and endpoint there is at most one possible
deterministic expression \eqref{eq:status-opposite-offset}, including
its coefficients, even if the permitted exclusions vary. Indeed two
candidate exclusion sets of size at most $B_*$ have a common surviving
generic label, since the label universe is larger than $2B_*$. At that
label every role/basis assignment has positive density conditional on
the signature. Subtracting the two claimed deterministic expressions
therefore forces their coefficient differences to be zero, and then
their offset functions agree almost everywhere on every $(l,q)$ space.
The same generic label can be used at every tag and diagonal bit.

Thus, for a fixed second-side array, each choice of $j$ supplies at most
four offsets for each first-side endpoint: the two multiples of $q$,
and the unique function \eqref{eq:status-opposite-offset} plus those
two multiples, if that function exists. Across both choices of $j$
there are at most $2\cdot4^2=32$ pairs $(f_1,f_2)$.

All events used here are measurable. There are finitely many types,
status vectors, coefficients, and exclusion sets, and generator
availability is a positivity test on the integral of jointly
measurable densities. By Fubini, choose a fixed second-side array for
which the first-side mass of $E_s$ exceeds $.04$. One of the at most
32 pairs of offset functions then gives
\eqref{eq:status-extracted-prediction} on first-side mass at least
$\delta_0=.04/32$, allowing its coefficients and exclusion sets to vary
with the first-side array. This is much larger than $10^{-6}$.

We now transfer these fixed signature functions to a prediction in the
finite problem. This step is necessary because the definition of
prediction asks for global functions of actual keys. For any
$\varepsilon>0$, approximate each fixed binary function $f_i$ by a
binary cylinder function $f_i^{(\varepsilon)}$, with disagreement
measure at most $\varepsilon$ on each of the finitely many $(l,q)$
reference spaces. Such approximation follows by approximating the
measurable event where $f_i=1$ in measure by the algebra of finite
cylinders. The total unary measure of every type is its reference
measure, so replacing $f_i$ increases every relation error in
\eqref{eq:status-extracted-prediction} by at most $\varepsilon$.

For fixed cylinder functions and a fixed choice of coefficient bits
and exclusions, all these error probabilities are finite linear
combinations of status-cylinder coordinates. Their maximum over the
remaining unary types is continuous. Minimizing over the finitely many
permitted coefficient/exclusion choices is also continuous, separately
on each discrete coordinate format. The strict event that this minimum
is below $2\varepsilon$ consequently has limiting probability at least
$\delta_0$. Weak convergence implies that its prelimit probability
exceeds $\delta_0/2$ for all sufficiently large indices. One may see
this directly by approximating the indicator of an open strict-error
event from below by continuous functions.

Choose decreasing $\varepsilon_m\to0$ and increasing prelimit indices
at which the last mass bound holds. At index $m$ use the two global
functions obtained by evaluating $f_i^{(\varepsilon_m)}$ on the actual
global key signature. Keep the orientations in the strict-error event
and choose one of their successful coefficient/exclusion records.
The kept mass is at least $\delta_0/2>10^{-6}$, and the maximum
relation error over all remaining unary types tends uniformly to zero.
Conditional-on-$q$ errors give the same conclusion for the unconditioned
flavor tests. These are precisely the requirements of
\Cref{def:prediction}. The no-prediction branch has only the first
effective spaces and only negligible trimming, so this mass margin
also gives a prediction before that trimming. This contradiction
proves the lemma.
\end{proof}

\subsection{The predicted dual relations}

In the predicted alternatives, exclude the stored prediction exceptions
as well as the mandatory support exclusions. Write the marked profiles
as $\lambda_{A,i}$ and $\lambda_{B,i}$. The constants $\alpha_i$ and
the functions $f_i$ are common by endpoint index across units. Put
$d_i=1+\alpha_i$. The prepared class $e_i=(d_i,[f_i])$ is as in
\eqref{eq:status-class}: functions are first identified almost
everywhere on all tag/bit signature spaces, and then quotiented by
the span of the one function $q$, with a common coefficient on those
spaces.

In a limiting table computation, use the notation
\begin{equation}
\label{eq:status-table-delta}
\mathsf u_{B,j}(\Delta_A)
 =\sum_{i=1}^2 (u_{B,j}U_{A,i})^*(\lambda_{A,i}),
\qquad
\mathsf u_{A,i}(\Delta_B)
 =\sum_{j=1}^2 (u_{A,i}U_{B,j})^*(\lambda_{B,j}).
\end{equation}
These are finite contractions in the recorded table coordinates. The
notation does not require retaining an ambient tensor $\Delta$ in the
compactification.

\begin{lemma}[Predicted scalar compatibility]
\label{lem:status-predicted}
For a predicted pair with a table satisfying the compatibility
conditions of \Cref{prop:status-preparation}, the system
\eqref{eq:status-local-targets}--\eqref{eq:status-coupled-targets}
is soluble. The assertion survives any further bounded label
exclusions leaving generic labels.
\end{lemma}

\begin{proof}
The exact limiting prediction is
\[
p_{A,i}=\boldsymbol\lambda_{A,i}\cdot\mathbf x_{A,i}
                 +\alpha_i h_{A,i}+f_i,
\qquad
p_{B,j}=\boldsymbol\lambda_{B,j}\cdot\mathbf x_{B,j}
                 +\alpha_j h_{B,j}+f_j,
\]
where bold $\boldsymbol\lambda$ denotes coordinates in the effective
basis. Substitute this into a dual relation
\eqref{eq:status-dual-relation}. At a fixed signature the two generic
status distributions have every pair of role/basis assignments
available. Comparing their coefficients forces
\begin{align}
\mathbf c_{ij}&=s_j\boldsymbol\lambda_{A,i},&
\mathbf c'_{ij}&=r_i\boldsymbol\lambda_{B,j},&
\beta_{ij}&=s_jd_i=r_id_j,
\label{eq:status-forced-coefficients}\\
s_jf_i+r_if_j&=\gamma_{ij}q.
\label{eq:status-forced-functions}
\end{align}
Thus
\begin{equation}
\label{eq:status-class-equations}
s_je_i=r_ie_j\qquad(i,j\in\{1,2\}).
\end{equation}
Equality of functions here is exactly the almost-everywhere equality
recorded in the prepared subsequential class.

If both $e_i$ vanish, \Cref{lem:empty-status} supplies the exact
identity $p_i=T(\lambda_i,\mathrm{atom})+h_i$ on all atoms. Hence on
each span $p_i+h_i=\boldsymbol\lambda_i\cdot\mathbf x_i$.
At the local targets, the first coupled sum for a fixed $j$ is
\[
\sum_i a(\lambda_{A,i})+
       \mathsf u_{B,j}(\Delta_A)
=c+\mathsf u_{B,j}(\Delta_A)=0
\]
by the empty-status compatibility condition. The reciprocal coupled
sums vanish in the same way. Local surjectivity therefore solves all
constraints in this case.

Suppose that some $e_i$ is nonzero. A nonzero solution of
\eqref{eq:status-class-equations} is possible only when all nonzero
$e_i$ are equal, and then both weight vectors are the indicator of
their support $S$. To check this without any assumption on the dimension
of the function space, first consider a support with one index, say
$e_1\ne0,e_2=0$. The off-diagonal equations force $r_2=s_2=0$,
and the diagonal gives $s_1=r_1=1$ for a nonzero solution. If both
$e_i$ are nonzero, the diagonal equations give $s_i=r_i$.
The off-diagonal equation $s_2e_1=s_1e_2$ either makes both weights
zero, or makes both one and $e_1=e_2$. These cases exhaust the
possibilities.

If the two nonzero classes are unequal, no dual obstruction exists.
Otherwise let $d$ be the common first bit on $S$. With
$s_i=r_i=\1_{i\in S}$, \eqref{eq:status-forced-functions} is symmetric
in $i,j$. Its diagonal is zero. Since the function $q$ is not zero
on the family of reference spaces, this gives
$\gamma_{ij}=\gamma_{ji}$ and $\gamma_{ii}=0$. The sum of all
four $q$-contributions to a target is therefore zero. The total role
contribution is
\[
\sum_{i,j}\beta_{ij}=|S|d\quad\text{in }\F.
\]

Evaluate the basis terms in
\eqref{eq:status-forced-coefficients} at their local targets.
The $a$-terms sum to $|S|c+|S|c=0$, because
\Cref{prop:status-preparation} fixed the same value
$c=\sum_i a(\lambda_i)$ in both units. The remaining terms are
\[
\sum_{j\in S}\mathsf u_{B,j}(\Delta_A)
 +\sum_{i\in S}\mathsf u_{A,i}(\Delta_B).
\]
By \eqref{eq:phase-compatibility}, this equals $|S|d$.
Together with the role contribution it is zero. Thus every dual
relation vanishes on the target, and \Cref{lem:status-duality} proves
solvability. The argument used only surviving generic labels and the
prediction on those labels, so additional exclusions of the indicated
kind do not change it.
\end{proof}

\begin{lemma}[Scalar solvability on positive mass]
\label{lem:status-span}
In each positive-mass alternative of
\Cref{prop:status-preparation}, the common experiment has positive mass
of passed injecting tables for which the scalar system is soluble,
robustly under the exclusions needed below. On the no-prediction branch
this holds whenever the total exclusions stay within $B_*$ per endpoint.
On a predicted branch it holds after the stored exceptions and any
further bounded exclusions leaving generic labels.
\end{lemma}

\begin{proof}
Without prediction, the probability of an obstructable base pair is
less than $.1$ by \Cref{lem:status-global-prediction}. This statement
uses the unflagged iid base law. In the full conditional-law experiment,
a passed injecting table exists on mass at least $.99$ by
\Cref{lem:limit-model}. Their intersection therefore has mass greater
than $.89$. On it no nonzero dual pattern exists for any exclusions
within the stated budget, so \Cref{lem:status-duality} solves the
scalar system for every such table. No independence of the two events
is being assumed.

With prediction, a passed compatible injecting table exists on positive
mass in the same experiment. The exact prediction and all almost-sure
unary properties continue to hold on that event, and
\Cref{lem:status-predicted} gives the assertion.
\end{proof}

\subsection{Stable spans and fresh short lists}

It remains to replace a vector in an edge span by a short list of actual
generators, using distinct labels across both lists at an endpoint.
The small dimension in \eqref{eq:status-full-projection} is decisive.

\begin{lemma}[Fresh lists of at most seven generators]
\label{lem:fresh-lists}
On any pair for which scalar solvability survives the bounded exclusions
specified in \Cref{lem:status-span}, further bounded exclusions can be made so that every scalar solution on the resulting
four edge spans can be realized by between one and seven available
generators per edge. The two endpoint labels of each generator are
retained with its certificate, and all labels used at an endpoint are
distinct across its two incident lists. Every realized matched
position has positive status overlap.
\end{lemma}

\begin{proof}
Begin with the mandatory support exclusions and, in the predicted
case, the stored exceptions. More generally, start with the initial
exceptions allowed by the scalar-solvability hypothesis. Whenever deleting at most 100 additional
labels at each endpoint can shrink any of the four edge spans, make
such a deletion. Each span continues to project onto its full base
space. Its dimension exceeds that base dimension by at most two, so
there are at most eight dimension-reducing steps in total. At the end,
every edge span is unchanged by deleting up to 100 further labels at
each endpoint. There remain generic labels, since
\eqref{eq:selector-size} makes the label universe much larger than all
these bounded exclusions. On the no-prediction branch the total budget,
including later fresh-label removals, is at most $B_0+900<B_*$.
On a predicted branch it is at most the stored exceptions plus this
quantity, still leaving many generic labels. Scalar solvability is
retained by the scalar-solvability hypothesis. Fix one solution on these stabilized
spans.

We prove a statement for one edge after any labels used on earlier
edges have also been reserved. Write its stabilized span as $V$, its
base space as $B$, and its projection as $\operatorname{pr}_B$. Let $G_L$ be the
generators avoiding the current reserved sets $L$. Their span is still
$V$. Let $H_L$ be the span of sums of one, two, or three members of
$G_L$ whose total base projection is zero and whose endpoint labels
are mutually disjoint within that word. We claim
\begin{equation}
\label{eq:status-short-kernel}
H_L=\ker\operatorname{pr}_B.
\end{equation}

To prove the claim, take two generator lifts of the same $b\in B$.
There is a third generic lift of $b$ at labels avoiding both of them
and $L$. Comparing each original lift to this third one gives an
allowed two-generator kernel word. Hence all lifts of $b$ have the
same class modulo $H_L$. Call that class $\ell(b)$. A generator
above base zero is itself an allowed one-generator kernel word, so
$\ell(0)=0$. Choose mutually label-disjoint generic lifts of
$b,c,b+c$. Their sum is an allowed three-generator kernel word,
and therefore $\ell(b)+\ell(c)=\ell(b+c)$. The generators span $V$,
so this additive lift identifies $V/H_L$ with $B$: the induced
projection is inverse to $\ell$. This proves
\eqref{eq:status-short-kernel}, including the case of zero base
projection.

For the desired vector in $V$, first choose a generic singleton lift
of its base projection. Reserve its labels. The remaining discrepancy
belongs to $\ker\operatorname{pr}_B$, whose dimension is at most two. On the remaining
labels \eqref{eq:status-short-kernel} still holds because stability has
preserved $V$. If the kernel is zero no correction is needed. If it is
one-dimensional, choose a nonzero short kernel word and use it if
needed. If it is two-dimensional, choose a nonzero short word, reserve
its at most three labels at each end, and apply
\eqref{eq:status-short-kernel} again. The new short words still span
the full kernel, so one lies outside the first word's direction.
These two words form a basis and the discrepancy is the sum of an
appropriate subset of them. This uses at most
$1+3+3=7$ generators. The initial lift is always retained, so the list
is nonempty; reserving labels before each choice ensures disjointness
even when two generators have the same numerical bit vector.

Carry out this procedure on the four edges in turn. Each endpoint is
incident to only two lists and uses at most fourteen labels in total.
All removals and auxiliary freshness comparisons in the procedure are
within the 100-label stability allowance. The generators keep their
certificates of positive overlap, proving the lemma.
\end{proof}

\subsection{Positive overlap on permitted pairs}

\begin{proposition}[Positive-mass overlap]
\label{prop:positive-overlap}
Consider two independent unit laws with actual leaf weights, the mixed
caps \eqref{eq:histogram-mixed-caps}, the final per-leaf image bounds,
and a permitted-pair indicator $J_n$ as in \Cref{sec:common-limit}.
Suppose their common model has positive mass of pairs with $J=1$
admitting a passed injecting table whose scalar system is soluble after
every bounded exclusion used by \Cref{lem:fresh-lists}. Then some fixed
numerical option has expected unnormalized, permitted-pair overlap
bounded below by a positive constant along a subsequence. Consequently
it satisfies \eqref{eq:overlap-criterion}, with the same permitted-pair
predicate. The two laws need not coincide. If the original laws
additionally satisfy the hypotheses of \Cref{prop:collision}, its
probability bound holds along the same subsequence, with any initial
restrictions of reciprocal mass $2^{o(N)}$ charged to that probability.
\end{proposition}

\begin{proof}
Work on the subsequence defining the common model in the hypothesis.
If the conclusion failed, finiteness of the option collection would
give \eqref{eq:limit-vanishing-overlap} along this same subsequence.
Apply \Cref{prop:limit-overlap} in this model.

By hypothesis, there is positive model mass of permitted pairs with a
table whose flags pass and whose scalar solution survives the needed
exclusions. Apply \Cref{lem:fresh-lists}. This gives between one and
seven matched generators per edge with fresh allowed labels. Their
sum vectors obey \eqref{eq:status-local-targets} and
\eqref{eq:status-coupled-targets}; equivalently, the lists obey all
scalar recipe equations. At each generator the unspecified tester
bits were only aggregated. Expanding the positive overlap integral
as a finite sum over the two full tester summaries yields at least
one pair of summaries for which the overlap is still positive.
Make that refinement independently at each position.

The scalar consistency argument in
\Cref{sec:realization}, in particular \eqref{eq:recipe-gradients},
now supplies a consistent numerical prescription of the required
within-endpoint binary bits from these summaries and $p$ bits. Together
with the already passed table, these data define a usable numerical
option. The only parameter filters added are precisely the unary
statuses and the prescribed binary conditions in
\eqref{eq:limit-success-format}. At every matched position the refined
status overlap is positive. This contradicts
\Cref{prop:limit-overlap} on a positive-mass event.

No pointwise uniform lower bound for these individual positive
integrals is needed. The options form a finite collection and every
product has finitely many nonnegative factors. A positive-mass event
on which some such product is positive gives positive expectation
for at least one option. Its lower bound is a constant on the chosen
subsequence, however small that constant may be. The preceding
contradiction therefore proves the stated subsequential lower bound.

For this probability consequence, assume additionally that the original
laws satisfy the hypotheses of \Cref{prop:collision}. The successful
densities are unnormalized with respect to the retained law, and
\Cref{prop:collision} explicitly charges any law
restriction of reciprocal mass $2^{o(N)}$. A positive constant, or
its product with such restriction costs, is $2^{-o(N)}$. The resulting
four-hole probability, including the permitted-pair predicate, has
exponent at most $k_{\max}+.03+o(1)$, whereas
\[
100g-(k_{\max}+.03)=44g+55.97>0.
\]
Thus these overlaps are on the scale required for
\Cref{thm:distribution-tests}.
\end{proof}


\subsection{The coloured unit distribution test}

\begin{proposition}[Different-colour conflicts]
\label{prop:coloured-distribution}
Let the early constant $\eta_{\mathrm{col}}>0$ satisfy
\[
 \eta_{\mathrm{col}}<\frac14c_* ,\qquad
 c_*:=\frac{\delta_*^2}{64},
\]
where $\delta_*$ is the retained-mass margin in
\eqref{eq:preparation-early-margins}. This choice depends on the original
endpoint bound $M$ and the early geometric constants, and precedes the
selector and channel dimensions. For every fixed finite palette and all
sufficiently large $n$, a unit law with both endpoint marginals bounded
by $M\pi$, pair law after forgetting colour bounded by $2^{3000gN}\pi^2$, and each
colour of mass at most $\eta_{\mathrm{col}}$ satisfies
\[
 \Pr\{A\text{ and }B\text{ have four cross holes and different colours}\}
 \ge 2^{-100gN},
\]
where $A,B$ are independent draws from that law. In particular,
\Cref{thm:distribution-tests}\textup{(i)} holds.
\end{proposition}

\begin{proof}
Suppose that the assertion fails along an unbounded sequence of $n$.
The laws and their randomized colour marks may vary with $n$; the
palette is fixed. Refine by colour and discard the exponentially light
colour classes as in \Cref{lem:peeling}. Peel each remaining conditional
class, but retain the original mass of every leaf. Their mixture is a
restriction of the original law with total mass $1-o(1)$ before the
fixed preparation restrictions. Hence its endpoint marginals remain
bounded by a fixed multiple of $M\pi$ and its primal marginals by the
same multiple of the original primal law. The conditioning of a small
colour may change a leaf's joint density bound; this is already included
in its absolute density and image budgets. It never changes the mixed
marginal constant used next.

Apply \Cref{prop:status-preparation} to this pooled law, carrying its
colour as a mark. In the predicted case this yields a retained law of
original mass at least $\delta_*$, and an event of passed compatible
injecting tables of conditional pair mass at least $1/64$. Its scalar
systems are soluble by \Cref{lem:status-predicted}, robustly under the
fresh-label exclusions. In the no-prediction case the retained mass is
also at least $\delta_*$ and \Cref{lem:status-span} supplies conditional
pair mass greater than $.89$. In both cases the successful scalar-table
event has mass at least $c_*$ in two draws from the original law.

If the original colour probabilities are $p_c$, their equal-colour
pair mass is
\[
 \sum_c p_c^2\le (\max_c p_c)\sum_c p_c
 \le\eta_{\mathrm{col}}.
\]
The equal-colour part of any subevent, including a prepared subevent,
has no larger original mass. Subtracting it therefore leaves at least
$c_*-\eta_{\mathrm{col}}>0$ of successful scalar-table pairs of different
colours. Equivalently, if preparation retained mass $\delta$, its
normalized equal-colour pair mass is at most
$\eta_{\mathrm{col}}/\delta^2$. This subtraction takes place before any
pigeonholing of a numerical table. Thus neither the number of tables nor
the palette size enters the early choice of $\eta_{\mathrm{col}}$.

Use the permitted-pair indicator
$J_n(L_A,L_B)=\mathbf1_{\{\operatorname{col}(L_A)
\ne\operatorname{col}(L_B)\}}$. Store this bit in the common model.
The preceding original-mass lower bound, together with the retained
prediction and injection assertions of \Cref{lem:limit-model}, gives
positive limiting mass on $J=1$ satisfying the hypothesis of
\Cref{prop:positive-overlap}. That proposition yields a fixed option
with positive subsequential overlap on different-colour pairs.
The filtered form of \Cref{prop:collision} now gives probability at
least $2^{-(k_{\max}+.03)N-o(N)}$ for a four-hole event on those same
pairs. Because $k_{\max}+.03<100g$, this is larger than
$2^{-100gN}$ for all sufficiently large indices, contradicting the
chosen violating sequence.

The overlap constant is allowed to depend on the subsequential model.
This causes no loss in the uniform eventual assertion: any sequence
violating it has a further subsequence with one fixed positive constant,
and the fixed exponential gap just used defeats that sequence.
\end{proof}
